\section{Inteligența Artificială în Modă (Fashion AI)}
\label{sec:ai_fashion}

Tranziția de la analiza textului la analiza imaginii a revoluționat modul în care algoritmii înțeleg moda. "Fashion AI" este un subdomeniu al viziunii artificiale (\textit{Computer Vision}) care aplică tehnici de Deep Learning pentru a rezolva probleme specifice industriei textile.

\subsection{Clasificarea și Detectarea Articolelor}
Pentru a construi o ținută, sistemul trebuie mai întâi să "vadă" și să înțeleagă ce haine deține utilizatorul.
\begin{itemize}
    \item \textbf{Arhitecturi CNN (Convolutional Neural Networks)}: Rețele precum ResNet50 \cite{he2016deep} sau EfficientNet \cite{tan2019efficientnet} sunt standardul de aur pentru clasificarea imaginilor. Acestea extrag ierarhic caracteristici, de la margini și texturi simple (straturile inferioare) până la forme complexe și obiecte (straturile superioare).
    \item \textbf{Detectarea Obiectelor (Object Detection)}: Algoritmi precum YOLO (You Only Look Once) \cite{redmon2016you} permit localizarea precisă a hainelor într-o imagine, esențială atunci când utilizatorul fotografiază o ținută completă.
\end{itemize}

\subsection{Segmentarea Semantică și Eliminarea Fundalului}
O problemă critică în digitalizarea garderobei este calitatea imaginilor realizate de utilizatori acasă (fundaluri aglomerate, iluminare slabă).
\textbf{U\textsuperscript{2}-Net} \cite{qin2020u2net} reprezintă stadiul actual al tehnologiei în detectarea obiectelor proeminente (\textit{Salient Object Detection}). Arhitectura sa, bazată pe structuri imbricate de tip U-Net, permite o segmentare pixel-perfect a hainei, eliminând fundalul cu o precizie superioară metodelor clasice de procesare a imaginii.

\subsection{Compatibilitatea Vizuală și Transferul de Stil}
Cercetările recente se concentrează pe modelarea matematică a compatibilității.
\begin{itemize}
    \item \textbf{Siamese Networks}: Rețele care învață o funcție de distanță între două obiecte, antrenate pe seturi de date de ținute (ex: Polyvore dataset), pentru a decide dacă două articole "merg împreună" \cite{vaswani2017attention}.
    \item \textbf{Color Harmony Models}: Abordări care extrag histogramele de culoare și le compară cu șabloane euristice de armonie (Itten, Moon & Spencer) pentru a valida estetic o combinație.
\end{itemize}

StyleGenAI adoptă o abordare pragmatică, utilizând CNN-uri pre-antrenate pentru extracția de atribute și un motor bazat pe reguli (\textit{rule-based engine}) pentru validarea cromatică, o metodă robustă în absența unor seturi de date masive etichetate.
